\section{Simulator Implementation}
\label{sec:simulator}

We implemented \GreenFaaS-sim, a single-process discrete-event
simulator in Python, alongside the scheduler. The design follows the
abstractions of \code{faas-sim} (function-level FaaS simulation) and
\code{vessim} (carbon co-simulation) \citep{faassim,vessim}, but is a
clean-room implementation rather than a fork --- written from scratch
to give us tight control over event ordering, warm-pool accounting,
and the extension points for trace loaders. The full implementation,
including all baseline schedulers, is approximately 1{,}700 lines of
substantive Python with no required dependencies beyond the standard
library.

\subsection{Event Model}

The simulator maintains a min-heap of timestamped events of four
kinds. \textsc{Arrival}: a new invocation enters the system; the
scheduler is consulted exactly once and the decision committed for
the lifetime of the invocation. \textsc{Start}: the chosen start time
of a previously-scheduled invocation; execution begins, energy and
carbon are charged. \textsc{Completion}: an invocation finishes; the
container becomes warm-and-idle for up to $T_w$ seconds.
\textsc{WarmExpire}: a warm container reaches its TTL without being
reused and is evicted; accumulated idle energy is finalised.

The decoupling of \textsc{Arrival} from \textsc{Start} is what makes
temporal deferral possible without special-casing in the event loop:
the scheduler returns a \code{start\_time} that may be in the
future, and the simulator simply enqueues the \textsc{Start} event at
that time. Reusing a warm container is signalled by the scheduler's
\code{use\_warm} flag and realised at the \textsc{Start} event by
decrementing the warm-pool count.

\subsection{Warm-Pool Accounting}

For each $(r, f)$ pair, the simulator tracks $W_{r,f}(t)$ and a queue
of pending \textsc{WarmExpire} events. A container becomes warm at
\textsc{Completion} with scheduled expiry $T_w$ seconds later. If
reused before expiry, the warm count is decremented and the
corresponding \textsc{WarmExpire} is matched at the \textsc{Start}
event; the container's lifetime
$(t_{\mathrm{end}} - t_{\mathrm{start}})$ is logged. If expiry fires
first, the \textsc{WarmExpire} event performs eviction and logs the
full $T_w$ lifetime.

At simulation end, total warm-container lifetime per $(r, f)$ pair
is multiplied by $P_i(f)$, $\mathrm{PUE}_r$, and the region's
\emph{mean} carbon intensity $\bar{C}_r$. As discussed in
\S\ref{sec:problem:cost}, this mean-intensity term is an
approximation of the exact interval integral
$\int W_{r,f}(t) C_r(t)\,dt$; it is exact when idle time is
distributed uniformly over the horizon and conservative otherwise.
The simulator therefore reports results under the mean-intensity
approximation justified in \S\ref{sec:problem:cost}, matching the
convention used in \GreenFaaS's per-slot scoring
(Algorithm~\ref{alg:scoreslots}, line 13).

\subsection{Extension Points}

Three abstractions cleanly separate the algorithm from the data
sources. \textbf{\code{CarbonModel}} provides
\code{intensity(region, t)} and \code{forecast(region, t, horizon, accuracy)};
the default implementation reads from per-region \code{CarbonTrace}
objects with arbitrary step sizes, and real LWA / ENTSO-E format CSVs
\citep{lwa} are loaded via \code{load\_carbon\_csv} into the same type.
The scheduler does not know whether the carbon data is synthetic or real.
\textbf{Workload generators} produce \code{List[Invocation]}; the
default emits a non-homogeneous Poisson stream with Zipf
popularity, and the real-trace loaders emit the same type from
published Azure schemas. \textbf{Schedulers} implement a single
method \code{schedule(invocation, state, carbon) -> ScheduleDecision};
all five schedulers in our evaluation conform to this interface and
are swappable.

\subsection{Validation}

The simulator's correctness rests on three checks. \emph{(1)
Conservation:} total reported execution energy equals
$\sum_i P_a(f(i)) \tau(i)$ within floating-point precision; total
warm-idle energy equals $\sum_{(r,f)} P_i(f) \cdot
\mathrm{lifetime}_{r,f}$. Verified by direct comparison with
hand-computed values on small synthetic workloads. \emph{(2) Schedule
determinism:} re-running a fixed scheduler on a fixed invocation
stream produces identical decisions; verified for all five
schedulers. \emph{(3) Lemma-implementation agreement:} the
closed-form Lemma~\ref{lemma:tradeoff} and the per-invocation carbon
computed by direct simulation agree at 56 grid points spanning
$r \in [1.1, 11.7]$ and $\lambda \in [10^{-4}, 1]$/s, with
\emph{zero} mismatches: the brute-force simulator confirms the
lemma's predicted regime boundary (warm-in-low-carbon versus
cold-in-high-carbon) at every grid point, and confirms the predicted
break-even rate $\lambda^*$ within the grid resolution
(\code{scripts/verify\_tradeoff.py}). The third check
in particular gives us confidence that the scheduler's lemma-driven
decisions and the simulator's energy accounting are mutually
consistent --- addressing a concern that arises because the lemma was derived analytically
and the simulator was implemented independently.
